%% 第24回編集委員会（物理第3巻）：組 波W8
%% p3-00.tex  波-0　序章　ゆれは止められる？　波は使える？（★1〜2。3段にしない）
%% 中身は計画書 tools/phys3-plan.md の B「波-0」。本文の章（波-1〜4・7〜13）を読んで書いた、この巻の地図。第2巻の序章 p2-00.tex と同じ形。
%% 序章なので、じぶんでとこう・章末問題はない（例題は7段の見本3つ）。「この章で導く中学・高校の公式」の欄は置かない
%%   （c＝fλ の持ち主は波-5 ②、弦の固有振動数は波-6 ③、v＝c/n は波-12 ③。第19回 6-1-2）。
%% 記号カードは置かない（周期 T・振動数 f と減衰係数 c は波-1、光速 c は波-12 がはじめに置いている。U-6）。
%% 記号（この章）：周期 T〔s〕、振動数 f〔Hz〕、波長 λ〔m〕、波の速さ c〔m/s〕（この章では c を「波の速さ」の意味だけで使い、
%%   光の速さも「光の速さ c」と書く）、屈折率 n、距離 d、時間 t、弦の長さ L（ふりこは出さない）。
%% 数値はすべて Python で計算した：checks/物理第3巻/波-0/calc_p3_00.py（出力 calc_p3_00_out.txt）。
%% 本文の章とそろえた数値（同じ題材は、持ち主の章の値）：
%%   ・雷 3.0 s・音速 340 m/s → 約1.0 km：波-9 ②・例題1 と同じ（中学理科とくらべる例題は 340 m/s。20 ℃の空気は 343 m/s：熱-4 ⑥・波-9）。
%%   ・30階のビルの固有周期 約3秒（固有振動数 約0.3 Hz）：波-1 の桁の感覚（工学ではこう使う）と同じ。目安「階数×0.1 秒」の出典は p3-01 の頭のコメント。
%%   ・ブランコ 2.0 m → 2.8 s：波-1 例題1。
%%   ・台北101 の重り（直径 5.5 m・660 t）：波-3（p3-03 の出典）。
%%   ・光ファイバのコア 1.48：波-12 ④・例題（p3-12 の出典）。光の速さ 3.00×10^8 m/s：波-12・波-8。
%%   ・ギターの A弦 110 Hz・0.65 m・c＝143 m/s：数学 9-1、波-6 ③（計画書 B 波-6 の例題1）。
%%   ・人に聞こえる 20 Hz〜20 kHz：波-11（p3-11 の出典）。超音波検査 2.5〜10 MHz：波-10（p3-10 の出典）。
%%   ・目に見える光 約380〜780 nm：波-13・熱-13（p3-13 の出典）。緑 550 nm → 5.45×10^14 Hz：波-13 ①。
%%   ・救急車 960 Hz：波-8 ⑤。超音波の検査のゼリー：波-7 ④。
%%   ・波-5・波-6 は別の組が書いている最中なので、計画書 B の節番号で指した（波-5 ①：波とは・P波とS波、②：c＝fλ、④：弦の波動方程式。
%%     波-6 ③：弦の固有振動）。
%% 出典（第24回の追記（11月5日）により、書くときに確かめた）
%%   ・地震の予知：気象庁「地震予知について」https://www.jma.go.jp/jma/kishou/know/faq/faq24.html（curl で本文を読んだ）
%%     「現在の科学的知見からは、そのような確度の高い地震の予測は難しいと考えられています」。
%%     地震調査研究推進本部の FAQ（https://jishin.go.jp/resource/faq/）も「一般的には地震予知は困難」。
%%   ・P波は秒速約7 km、S波は秒速約4 km：気象庁「緊急地震速報のしくみ」https://www.jma.go.jp/jma/kishou/know/jishin/eew/shikumi/shikumi.html
%%     と、気象庁大阪管区気象台「地震一口メモ No.239」（PDF）を curl で読んだ。
%%   ・長周期地震動（周期が長い大きなゆれで、高層ビルが大きく長くゆれる。数百 km はなれた所でも）：気象庁「長周期地震動について」
%%     https://www.data.jma.go.jp/eqev/data/choshuki/choshuki_eq1.html（curl で本文を読んだ）。
%%   ・中波（AM）放送 531〜1602 kHz（9 kHz おき）：総務省「基幹放送用周波数使用計画」の検索の抜粋、総務省関東総合通信局（526.5〜1606.5 kHz の帯）。
%%     FM 放送 76.1〜94.9 MHz：総務省東海総合通信局の報道資料（2024年1月9日）https://www.soumu.go.jp/soutsu/tokai/kohosiryo/2024/0109.html（curl で読んだ）。
%%   ・無線 LAN の 2.4 GHz 帯（2400〜2483.5 MHz）：総務省の周波数表（検索の抜粋。波-14 の頭のコメントと同じ）。
%%   ・〔要確認〕なし。
\makeatletter
% 会話の1行  \rxatalk{y}{sora|mina|delta}{名前}{せりふ}（p2-00 の \qxatalk と同じ形）
\newcommand\rxatalk[4]{%
  \bkChar{#2}{(7,#1-6.5)}{12mm}
  \node[rectangle callout, callout absolute pointer={(15.5,#1)}, anchor=west, draw=bkTheme!80!black, line width=0.8pt,
    fill=white, rounded corners=2.5mm, inner xsep=3mm, inner ysep=1.8mm, text width=103mm, align=left,
    font=\small] at (21,#1) {#4};
  \node[anchor=north, font=\sffamily\footnotesize, text=bkGray, inner sep=0pt] at (7,#1-7.5) {#3};}
% 高いビル  \rxabldg{(左下)}{傾き(mm)}{色}  上ほど大きくずれる形（1次のゆれ）
\newcommand\rxabldg[4][]{\begin{scope}[shift={#2}]
  \filldraw[fill=#4!14, draw=#4!80!black, line width=0.7pt, #1] (0,0) -- (10,0) .. controls (10,12) and ({10+#3*0.5},24) .. ({10+#3},34)
     -- ({#3},34) .. controls ({#3*0.5},24) and (0,12) .. (0,0) -- cycle;
  \foreach \yy in {5,10,...,30} {\draw[#4!60!black, line width=0.3pt, #1] ({#3*(\yy/34)*(\yy/34)},\yy) -- ({10+#3*(\yy/34)*(\yy/34)},\yy);}
\end{scope}}
\makeatother
\tikzset{rxa lab/.style={font=\sffamily\footnotesize, inner sep=1pt, text=bkInk},
  rxa arr/.style={-{Stealth[length=2.6mm, width=2.2mm]}, line width=1.1pt},
  rxa dim/.style={bkGray!80!black, line width=0.5pt, {Stealth[length=1.8mm]}-{Stealth[length=1.8mm]}}}

\phchapter{0}{ゆれは止められる？\\波は使える？}{地震のゆれ、楽器の音、光。どこが同じ？}{\bklevelrange{1}{2}}{ゆれは止められる？　波は使える？}
\bkmichishirube{2-10,4-8,力-2,力-4,熱-4}{波-1,波-2,波-5,波-9,波-12}

\begin{goal}
\begin{itemize}[label={\color{bkTheme}●}, leftmargin=1.6\zw, itemsep=0pt, topsep=1pt]
\item \bklevel{1}\enspace ゆれの周期と振動数、波の波長と速さを、身近な例で説明できる。
\item \bklevel{2}\enspace 何を無視したかを書いてから、波が届くまでの時間や距離を、7つの段で求められる。
\item \bklevel{1}\enspace 身近な「なぜ？」の答えが、この巻のどの章にあるかを探せる。
\end{itemize}
\end{goal}

\begin{hatena}
\noindent\begin{minipage}[c]{0.54\linewidth}\raggedright
地震が来ると、高いビルはゆっくり大きくゆれます。ギターの弦をはじくと、弦がふるえて音が出ます。スマホの動画は、光ファイバの中を進む光にのって、遠くから届きます。\par\vspace{3pt}
ビルのゆれ、楽器の音、光。見た目はまったくちがいますが、物理の目で見ると、同じところがあります。どこが同じなのでしょう。
\end{minipage}\hfill
\begin{minipage}[c]{0.43\linewidth}\centering
\begin{tikzpicture}[x=0.9mm, y=0.9mm]
  \fill[bkGray!25] (-2,-2) rectangle (58,0);
  \draw[bkInk, line width=0.6pt] (-2,0) -- (58,0);
  % ビル
  \rxabldg{(0,0)}{4}{bkBlue}
  % ギター（胴と弦）
  \filldraw[fill=bkOrange!30, draw=bkOrange!80!black, line width=0.7pt] (27,8) circle (6.5);
  \fill[white] (27,8) circle (2);
  \draw[bkOrange!80!black, line width=0.7pt] (27,8) circle (2);
  \filldraw[fill=bkOrange!50, draw=bkOrange!80!black, line width=0.6pt] (25.8,13) rectangle (28.2,30);
  \draw[bkInk, line width=0.6pt] (27,3) .. controls (28.5,14) and (28.5,22) .. (27,30);
  % 光ファイバ
  \draw[bkInk, line width=0.7pt] (40,6) .. controls (48,6) and (50,20) .. (57,22);
  \draw[bkInk, line width=0.7pt] (40,9) .. controls (47,9) and (49,23) .. (57,25);
  \fill[bkRed!70] (38,7.5) circle (1.4);
  \node[font=\sffamily\bfseries\Large, text=phHatena!80!black] at (17,26) {?};
  \node[font=\sffamily\bfseries\Large, text=phHatena!80!black] at (40,27) {?};
\end{tikzpicture}
\end{minipage}
\end{hatena}

この本は、身の周りの「ゆれ」と「波」の「なぜ？」に、物理で答える本です。序章では、地震のゆれ・楽器の音・光を例にして、この本で何を学ぶか、どの順に学ぶかを見ていきます。詳しい計算は、それぞれの章でします。いまはわからない式やことばが出てきても、だいじょうぶです。

\begin{yattemiyou}[波が運ぶもの、運ばないものを見よう]
\begin{enumerate}[label={\bkcircled[bkGreen]{\arabic*}}, leftmargin=2.4\zw, itemsep=1pt, topsep=2pt]
\item 洗面器かおふろに水を張り、水面に小さく切った紙を1枚うかべる。
\item 紙から\qty{20}{cm}ほど離れた所を、指で1秒に1回くらいの速さで、軽くたたき続ける。
\item 輪の形の波が、紙の所まで進んでいくのを見る。紙は、波といっしょに遠くへ運ばれていくだろうか。
\item たたく速さを1秒に2回にすると、輪と輪の間かくがどう変わるかを見る。
\end{enumerate}
\end{yattemiyou}

\noindent 紙は、その場で上下にゆれるだけで、遠くへは運ばれません。進んでいくのは、水そのものではなく、「ゆれている」という\textbf{形}です。たたく速さを2倍にすると、輪と輪の間かくはおよそ半分になります。この2つの見つけたことは、ビルのゆれにも、音にも、光にもあてはまります。この章の②で、そのわけを見ます。

\begin{omoidasou}
\bkrecall{4-8}{回る点の影は、$\sin$ の波の形になる（4-8⑥）。1周の時間が周期 $T$、1秒に回る回数が $f=1/T$（4-8⑨）。}
\bkrecall{2-10}{単位の前の k（キロ）は1000倍、m（ミリ）は $\tfrac{1}{1000}$、{\textmu}（マイクロ）は100万分の1（2-10③）。}
\bkrecall{力-2}{物の動きは、運動方程式 $ma=F$ で決まる。ゆれも波も、この式から始まる。}
\bkrecall{熱-4}{空気の中の音の速さは、\qty{20}{\celsius}で約\qty{343}{m/s}（⑥）。中学理科では「約\qty{340}{m/s}」と習う。}
\end{omoidasou}

%% ================================================================
\Needspace{16\baselineskip}
\bksection{hon}{この巻の地図}

\bksubsection{1}{3つの場面：ビル・ギター・光ファイバ}
\leavevmode\label{ph:波-0:sec1}はてな？の3つの場面を、もう少し詳しく見てみましょう（図\ref{ph:波-0:fig-scene}）。
\begin{bkfigure}
\begin{tikzpicture}[x=1mm, y=1mm]
  % (a) ビル
  \fill[bkGray!25] (0,-2) rectangle (36,0);
  \draw[bkInk, line width=0.6pt] (0,0) -- (36,0);
    \rxabldg[dashed, fill=none]{(12,0)}{-6}{bkBlue}
  \rxabldg{(12,0)}{6}{bkBlue}
  \draw[rxa arr, bkBlue!80!black] (6,30) -- (2,30);
  \draw[rxa arr, bkBlue!80!black] (28,30) -- (32,30);
  \node[rxa lab, align=center] at (18,-7) {(a)\enspace 地震のあと、ビルは\\1往復約3秒でゆれ続ける};
  % (b) ギター
  \draw[bkInk, line width=0.9pt] (44,26) -- (44,6) (84,26) -- (84,6);
  \draw[bkInk, line width=0.7pt] (44,16) .. controls (57,26) and (71,26) .. (84,16);
  \draw[bkGray, line width=0.6pt, dashed] (44,16) .. controls (57,6) and (71,6) .. (84,16);
  \fill[bkInk] (44,16) circle (0.9) (84,16) circle (0.9);
  \node[rxa lab, anchor=south] at (64,26) {弦の長さ \qty{0.65}{m}};
  \node[rxa lab, align=center] at (64,-7) {(b)\enspace 両はしが止まった弦が\\1秒に110回ふるえる};
  % (c) 光ファイバ
  \fill[bkBlue!8] (92,10) rectangle (132,22);
  \fill[bkBlue!22] (92,13) rectangle (132,19);
  \draw[bkInk, line width=0.6pt] (92,10) -- (132,10) (92,22) -- (132,22);
  \draw[bkBlue!70!black, line width=0.4pt] (92,13) -- (132,13) (92,19) -- (132,19);
  \draw[bkRed, line width=0.9pt] (92,16) -- (98,19) -- (110,13) -- (122,19) -- (132,14);
  \node[rxa lab, anchor=south] at (112,23) {ガラスの細い線};
  \node[rxa lab, align=center] at (112,-7) {(c)\enspace 光が、はね返りながら\\1秒に約20万 km 進む};
\end{tikzpicture}
\bkfigcaption{この巻の3つの場面。(a) ゆれる（実線と破線は、ゆれのはしの形）、(b) 弦のふるえが音になる、(c) 光が情報を運ぶ。}\label{ph:波-0:fig-scene}
\end{bkfigure}
\begin{enumerate}[label={\bkcircled[bkBlue]{\arabic*}}, leftmargin=2.4\zw, itemsep=2pt, topsep=2pt]
\item \textbf{ビルのゆれ}　30階建てくらいのビルは、地震のあと、1往復に約3秒かけて、ゆっくりゆれ続けます。ゆれが大きくなりすぎないように、ビルにはさまざまな「ゆれを止めるしくみ」がつけてあります。台湾の台北101 というビルの上のほうの階には、直径\qty{5.5}{m}、重さ660 t の鋼の球がつり下げてあり、ビルのゆれをおさえています（波-3）。
\item \textbf{ギターの音}　ギターのいちばん太い弦の1つ（A弦）は、長さ\qty{0.65}{m}で、1秒に110回ふるえます。弦のふるえが空気をゆらし、そのゆれが空気の中を伝わって、耳に届きます。指で弦をおさえて短くすると、音が高くなります（波-6）。
\item \textbf{光ファイバ}　髪の毛ほどの細いガラスの線の中を、光がはね返りながら進みます。インターネットの動画や声のデータの多くは、光の点めつとして、この線の中を運ばれています（波-12・波-14）。
\end{enumerate}
\noindent ①は、1か所で行ったり来たりをくり返す\textbf{ゆれ}（\textbf{振動}）です。\bkidx{しんどう}{振動}②と③は、ゆれが遠くへ伝わっていく\textbf{波}です。\bkidx{なみ}{波}②の弦のふるえはゆれで、それが空気を伝わる音は波です。ゆれと波は、いつもいっしょに出てきます。

\bksubsection{2}{ゆれと波は同じ形：くり返す・伝わる}
\leavevmode\label{ph:波-0:sec2}ばねにつるしたおもりを、少し引いてはなすと、上下にゆれます。そのおもりの高さを、時間を横軸にしてグラフにすると、4-8 で見た $\sin$ の形になります（図\ref{ph:波-0:fig-same}の左）。1往復にかかる時間を\textbf{周期} $T$〔s〕、1秒あたりの往復の回数を\textbf{振動数} $f$〔Hz〕といいます。\bkidx{しゅうき}{周期}\bkidx{しんどうすう}{振動数}1往復に $T$ 秒かかるなら、1秒では $1/T$ 回なので、$f=1/T$ です。Hz（ヘルツ）は「1秒に何回」の単位です。

長いひもの片はしを、同じリズムで上下にゆらし続けると、ひもに $\sin$ の形の波ができて、その形が進んでいきます（図\ref{ph:波-0:fig-same}の右）。ある時刻に、ひも全体の形を写真にとると、同じ形が一定の長さごとにくり返しています。このくり返しの長さを\textbf{波長} $\lambda$（ラムダ）〔m〕といいます。\bkidx{はちょう}{波長}ひもの1つの点に目をつけると、その点は、おもりと同じように上下にゆれているだけです。やってみようの紙と同じです。
\begin{bkfigure}
\begin{tikzpicture}
\begin{axis}[bkplot, at={(0mm,0mm)}, width=62mm, height=40mm, xmin=0, xmax=2.3, ymin=-1.5, ymax=1.6, xtick={1,2}, xticklabels={$T$,$2T$}, ytick={-1,1}, yticklabels={$-A$,$A$}, xlabel={時間 $t$}, ylabel={おもりの位置}]
  \addplot[bkcurve, domain=0:2.25, samples=100] {cos(360*x)};
  \draw[rxa dim] (axis cs:0,1.3) -- (axis cs:1,1.3);
  \node[rxa lab, anchor=south] at (axis cs:0.5,1.32) {周期 $T$};
\end{axis}
\begin{axis}[bkplot, at={(70mm,0mm)}, width=62mm, height=40mm, xmin=0, xmax=2.3, ymin=-1.5, ymax=1.6, xtick=\empty, ytick={-1,1}, yticklabels={$-A$,$A$}, xlabel={位置 $x$}, ylabel={ひもの高さ}]
  \addplot[bkcurve, domain=0:2.25, samples=100] {sin(360*x)};
  \addplot[bkthin, dashed, domain=0:2.25, samples=100] {sin(360*(x-0.25))};
  \draw[rxa dim] (axis cs:0.25,1.3) -- (axis cs:1.25,1.3);
  \node[rxa lab, anchor=south] at (axis cs:0.75,1.32) {波長 $\lambda$};
  \draw[rxa arr, bkOrange!80!black] (axis cs:1.55,1.35) -- (axis cs:1.95,1.35);
  \node[rxa lab, anchor=west] at (axis cs:1.97,1.35) {$c$};
\end{axis}
\end{tikzpicture}
\bkfigcaption{左：ばねのおもりのゆれ（時間のグラフ）。右：ひもを進む波（位置のグラフ。実線は今、破線は $\tfrac14T$ 後の形）。どちらも同じ $\sin$ の形で、横軸の意味だけがちがう。}\label{ph:波-0:fig-same}
\end{bkfigure}
\noindent ゆれの大きさ（中心からいちばん離れた所までの距離）を\textbf{振幅} $A$ といいます。\bkidx{しんぷく}{振幅}ゆれは時間の $\sin$、波は時間と場所の両方の $\sin$ で書けます。これが、この巻の背骨です。

\Needspace{14\baselineskip}
\phboxhead{波の速さ $c$ と、振動数・波長の関係}
波の形は、ひもの点が1往復する間（周期 $T$ の間）に、ちょうど1波長 $\lambda$ だけ進みます。図\ref{ph:波-0:fig-same}の右で、山が1つ前の山の位置まで進むと、形はもとと同じに見えるからです。速さは「進んだ距離÷かかった時間」なので、\textbf{波の速さ} $c$〔m/s〕は、次のようになります。\bkidx{なみのはやさ}{波の速さ}
\begin{phflow}[ゆれから波の速さへ]
\phflowrow{波の形は、1周期 $T$ の間に\\1波長 $\lambda$ だけ進む}{形を変えずに\\一定の速さで進む}{$c=\dfrac{\lambda}{T}$}
\phflowrow{$c=\dfrac{\lambda}{T}$}{振動数 $f=1/T$\\を入れる}{$c=f\lambda$}
\end{phflow}
\pheq{c=f\lambda}{m/s ＝ (1/s) × m}
\label{ph:波-0:cfl}高校物理基礎で習う式です。この巻では、波の速さを $c$ と書きます（高校の教科書の $v$ や $V$ と同じ）。波動方程式という一般の式から、この式を導き直すのは波-5 ②です。
\begin{phyomi}[$c=f\lambda$]
1秒に $f$ 回ゆれると、1秒に $f$ 個の波（山と谷の組）が送り出される。1個の長さが $\lambda$ なので、1秒に進む長さは $f\lambda$。
\end{phyomi}
\begin{phmijika}[ピアノの「ラ」の音の波長]
楽器の音合わせに使う\qty{440}{Hz}の「ラ」の音は、\qty{20}{\celsius}の空気（$c=\qty{343}{m/s}$）の中で、波長が $343\div440\approx\qty{0.78}{m}$ です。1オクターブ低い\qty{220}{Hz}の「ラ」は、波長が2倍の\qty{1.56}{m}です。空気の中では、高い音も低い音も同じ速さで進むので、振動数が半分なら波長が2倍になります。やってみようで、たたく速さを2倍にすると輪の間かくが半分になったのも、同じわけです。
\end{phmijika}

\begin{phbaisuru}[半分]{弦の長さ $L$}
両はしが止まった弦では、弦の長さ $L$ の2倍が、いちばん低い音の波長になります（くわしくは波-6 ③）。弦を伝わる波の速さ $c$ が同じなら、$c=f\lambda$ から
\[f=\frac{c}{2L}\]
です。A弦（$L=\qty{0.65}{m}$、$c=\qty{143}{m/s}$）は $143\div1.30=\qty{110}{Hz}$。まん中をおさえて長さを半分にすると、振動数は2倍の\qty{220}{Hz}になり、1オクターブ高い音になります。
\begin{bkcenter}
\begin{tikzpicture}
\begin{axis}[bkplot, width=100mm, height=42mm, xmin=0.2, xmax=0.72, ymin=0, ymax=380, xtick={0.325,0.433,0.65}, xticklabels={0.325,0.433,0.65}, ytick={110,220,330}, xlabel={弦の長さ $L$（m）}, ylabel={振動数 $f$（Hz）}]
  \addplot[bkcurve, domain=0.21:0.71, samples=80] {143/(2*x)};
  \fill[bkRed] (axis cs:0.65,110) circle (2pt) (axis cs:0.325,220) circle (2pt) (axis cs:0.4333,165) circle (2pt);
  \node[rxa lab, anchor=north] at (axis cs:0.65,98) {もとの長さ：\qty{110}{Hz}};
  \node[rxa lab, anchor=south west] at (axis cs:0.335,232) {半分：\qty{220}{Hz}（2倍）};
  \node[rxa lab, anchor=south west] at (axis cs:0.445,178) {$\tfrac23$：\qty{165}{Hz}（1.5倍）};
\end{axis}
\end{tikzpicture}
\bkfigcaption{弦の長さと振動数（A弦、$c=\qty{143}{m/s}$）。長さに反比例するので、短くするほど急に高くなる。}
\end{bkcenter}
\end{phbaisuru}

%% ================================================================
\bksubsection{3}{ゆれは止められる？　波は使える？}
\leavevmode\label{ph:波-0:sec3}この序章の題の2つの問いに、この巻がどう答えるかを、先に見ておきましょう。

\phboxhead{ゆれは止められる？}
ゆれを0にすることはできません。地面がゆれれば、その上の物も、少しはゆれます。でも、ゆれを\textbf{小さくする}ことはできます。この巻では、次の3つの方法を学びます。
\begin{enumerate}[label={\bkcircled[bkBlue]{\arabic*}}, leftmargin=2.4\zw, itemsep=2pt, topsep=2pt]
\item \textbf{リズムを外す}　ゆすられるリズムと、物の固有振動数をずらすと、共振が起きず、ゆれは大きくなりません。洗濯機や機械を、やわらかいゴムの足で支えるのも、この考えです（波-2）。
\item \textbf{ゆれを横取りする}　台北101 の鋼の球のように、ビルと同じリズムでゆれる重りをつけると、重りがビルのゆれのエネルギーを受けとって、ビルのゆれが小さくなります（波-3）。
\item \textbf{測って打ち消す}　ゆれを測り、そのゆれを打ち消すように、機械で力を加えます。エアコンが部屋の温度を測って調節するのと同じ、フィードバックの考えです（波-4）。
\end{enumerate}

\phboxhead{波は使える？}
波は、遠くへ物を運ばずに、エネルギーと情報を運びます（波-5・波-7）。だから、波は「見えない所を調べる」「遠くへ知らせる」道具になります。
\begin{itemize}[leftmargin=1.2\zw, itemsep=2pt, topsep=2pt, label={\color{bkTheme}●}]
\item \textbf{調べる}：超音波のこだまで、おなかの中の赤ちゃんや、鉄の中の傷を見る（波-7・波-10）。救急車の音の高さが変わるのと同じしくみで、ボールの速さを電波で測る（波-8）。
\item \textbf{聞く・数にする}：音の大きさを dB で測る（波-9）。音をマイクで数の列にして、スマホに記録する（波-11）。
\item \textbf{見る・送る}：レンズで光を曲げて像を作る（波-12）。光の波の重なりで、髪の毛の太さの1万分の1より細かい長さを測る（波-13）。電波や光に声や動画をのせて送る（波-14）。
\end{itemize}

\begin{phmada}[地震がいつ来るかは、まだ答えられない]
地震のゆれがどう伝わり、建物がどうゆれるかは、この巻の物理でかなりの所まで計算できます。地震のゆれには、先に着く速い波（\textbf{P波}、地面の中を秒速約7 km）と、あとから着く大きなゆれの波（\textbf{S波}、秒速約4 km）があります（波-5 ①）。\bkidx{ぴーは}{P波}\bkidx{えすは}{S波}緊急地震速報は、震源の近くで先に P波をとらえて、S波が届く前に知らせるしくみです。でも、これは地震が起きた\textbf{あとに}知らせるもので、地震が起きる前に知らせるものではありません。「いつ、どこで、どれくらいの大きさの地震が起きるか」を、前もって精度よく知らせることは、気象庁も「現在の科学的知見からは難しい」と説明しています。地面の下で岩がこわれるしくみの研究は続いていますが、地震の予知は、まだできていません。
\end{phmada}
\begin{motto}[2]
\phboxhead{2つの波のおくれで、震源までの距離がわかる}
地震のとき、はじめに小さくカタカタとゆれ（P波）、少しおくれて大きくユサユサとゆれます（S波）。2つの波は同じ所から同時に出ますが、速さがちがうので、遠い所ほど、着く時刻の差が大きくなります。震源までの距離を $d$〔km〕とすると、P波は $d/7$ 秒、S波は $d/4$ 秒で着くので、着く時刻の差は
\[\frac{d}{4}-\frac{d}{7}=\frac{3d}{28}\ \text{〔s〕}\]
です。差が\qty{10}{s}なら、$d=10\times28\div3\approx\qty{93}{km}$。差が1秒長くなるごとに、震源は約\qty{9}{km}遠くなります（図\ref{ph:波-0:fig-ps}）。中学理科で習う「初期微動の長さは、震源までの距離に比例する」は、このことです。
\begin{bkcenter}
\begin{tikzpicture}
\begin{axis}[bkplot, width=84mm, height=44mm, xmin=0, xmax=210, ymin=0, ymax=58, xtick={0,50,100,150,200}, ytick={10,20,30,40,50}, xlabel={震源までの距離（km）}, ylabel={着くまでの時間（s）}]
  \addplot[bkcurve, domain=0:200] {x/7};
  \addplot[bkcurve2, dashed, domain=0:200] {x/4};
  \draw[rxa dim] (axis cs:93.3,13.3) -- (axis cs:93.3,23.3);
  \node[rxa lab, anchor=north] at (axis cs:100,10.5) {差 10\,s};
  \node[rxa lab, anchor=west] at (axis cs:150,15) {P波（実線）};
  \node[rxa lab, anchor=east] at (axis cs:165,47) {S波（破線）};
\end{axis}
\end{tikzpicture}
\bkfigcaption{震源からの距離と、P波・S波が着くまでの時間（P波 秒速約7 km、S波 秒速約4 km）。2本の直線の間（たての差）が、初期微動の長さ。}\label{ph:波-0:fig-ps}
\end{bkcenter}
\end{motto}

%% ================================================================
\bksubsection{4}{例題の解き方：7つの段}
\leavevmode\label{ph:波-0:sec4}この本の例題は、いつも同じ7つの段で解きます。①場面、②何を無視するか（理想化）、③図と記号、④使う式を選ぶ理由、⑤式を立てる、⑥解く、⑦答えを確かめる（単位・大きさ・極端な場合）です。数学の本の例題よりも段が多いのは、式を立てる前の「モデルを作る」部分と、答えを出したあとの「確かめる」部分を、ていねいに書くためです。何を無視したかは、式を立てる前に、次のような「モデル化の札」にも書きます。
\begin{phmodel}[雷の光と音]
\phignore{光が届くまでの時間。風。気温による音の速さのちがい。}
\phassume{音は、雷の光った所から、まっすぐ一定の速さで進む。}
\phlimit{雷がとても近いとき（稲妻は何 km も長いので、どこを「雷の所」とするかで変わる）。}
\end{phmodel}
\bkidx{もでる}{モデル}このように、知りたいことに関係の小さいものを思い切って無視して作った、簡単にした考え方を\textbf{モデル}といいます。見本を3つ見ておきましょう。

\begin{phreidai}[雷までの距離]\label{ph:波-0:ex1}
稲妻が光ってから3.0秒後に、雷の音が聞こえました。雷までの距離は約何 km ですか。音の速さは\qty{340}{m/s}とします。
\phbamen
雷の光と音は、同じ所から同時に出ます。光はすぐ届き、音はおくれて届きます。
\phmushi
モデル化の札のとおり、光が届くまでの時間と、風を無視します。音の速さは一定とします。
\phzu
雷までの距離 $d$〔m〕、音の速さ $c=\qty{340}{m/s}$、おくれ $t=\qty{3.0}{s}$。
\phriyu
音は一定の速さで進むので、「距離＝速さ×時間」が使えます。
\phtateru
\[d=c\,t\]
\phtoku
\[d=340\times3.0=\qty{1020}{m}\approx\qty{1.0}{km}\]
\phtashikame
単位：m/s × s ＝ m。\\
大きさ：光が\qty{1020}{m}進むのにかかる時間は $1020\div(3.00\times10^8)\approx\qty{3.4e-6}{s}$ で、3秒にくらべてずっと小さいので、無視してよかったことがわかります。\\
極端な場合：光ってすぐ音がしたら（$t$ がほぼ0）、雷はすぐ近くです。音の速さを\qty{343}{m/s}にしても、答えは1\,\%ほどしか変わりません。
\kotae 約\qty{1.0}{km}
\end{phreidai}
\noindent 「秒数を3でわると、およその km」と覚えておくと便利です。雷の音がなぜおくれるのか、音の速さは何で決まるのかは、波-9 ②で詳しく調べます（波-9 の例題1も同じ数です）。

\begin{phreidai}[30階建てのビルは、何秒で1往復ゆれる？]\label{ph:波-0:ex2}
高いビルが横に1往復ゆれる時間（いちばんゆれやすい周期）は、およそ「階数×0.1秒」であることが知られています。30階建てのビルの周期と振動数を求めましょう。
\phbamen
地震のあと、ビルが自分のリズムで、横にゆれ続けています。
\phmushi
ビルの形・材料・地面のちがいを無視して、目安の式「階数×0.1秒」を使います。目安なので、本当のビルでは、この0.6倍から1.5倍くらいまでの幅があります。
\phzu
階数 $N=30$、周期 $T$〔s〕、振動数 $f$〔Hz〕。
\phriyu
周期は目安の式から、振動数は $f=1/T$（②）から出せます。
\phtateru
\[T=0.1\times N,\qquad f=\frac{1}{T}\]
\phtoku
\[T=0.1\times30=\qty{3.0}{s},\qquad f=\frac{1}{3.0}\approx\qty{0.33}{Hz}\]
\phtashikame
単位：s の逆数で Hz（1秒に何回）。\\
大きさ：1分間に約20往復です。長さ\qty{2.0}{m}のブランコの1往復（約\qty{2.8}{s}、波-1 の例題1）と同じくらいの、ゆっくりしたゆれです。\\
極端な場合：10階建てなら約1秒、60階建てなら約6秒。高いビルほど、ゆっくりゆれます。
\kotae 周期 約\qty{3.0}{s}、振動数 約\qty{0.33}{Hz}
\end{phreidai}
\noindent 建物が自分のリズムでゆれるときの振動数を\textbf{固有振動数}といいます。\bkidx{こゆうしんどうすう}{固有振動数}大きな地震では、周期が数秒の、ゆっくりした大きなゆれ（\textbf{長周期地震動}）が起きることがあり、震源から数百 km 離れた所でも、高いビルが大きく長くゆれることがあります。\bkidx{ちょうしゅうきじしんどう}{長周期地震動}ビルの周期と、地面のゆれの周期が近いと、ゆれがどんどん大きくなるからです。これを\textbf{共振}といい、波-2 で学びます。

\begin{phreidai}[光ファイバで 400 km 送る時間]\label{ph:波-0:ex3}
東京から大阪まで、長さ\qty{400}{km}の光ファイバで、声のデータを光にのせて送ります。光ファイバのガラスの屈折率は1.48です。光が着くまでの時間を求めましょう。真空の光の速さは\qty{3.00e8}{m/s}です。
\phbamen
光の点めつが、ガラスの細い線の中を進んでいきます。
\phmushi
光が線の中ではね返りながら進むジグザグの道と、送る機械・受ける機械の中でかかる時間を無視します。
\phzu
長さ $L=\qty{400}{km}=\qty{4.00e5}{m}$、真空の光の速さ $c=\qty{3.00e8}{m/s}$、屈折率 $n=1.48$、ガラスの中の光の速さ $v$、時間 $t$。
\phriyu
ガラスの中では、光は真空の $1/n$ の速さになります（$v=c/n$。波-12 ③）。速さがわかれば「時間＝距離÷速さ」です。
\phtateru
\[v=\frac{c}{n},\qquad t=\frac{L}{v}=\frac{nL}{c}\]
\phtoku
\[v=\frac{3.00\times10^8}{1.48}\approx\qty{2.03e8}{m/s}\]
\[t=\frac{4.00\times10^5}{2.03\times10^8}\approx\qty{1.97e-3}{s}\approx\qty{2.0}{ms}\]
\phtashikame
単位：m ÷（m/s）＝ s。\\
大きさ：光は、ガラスの中でも1秒に約20万 km 進みます。電波で空気の中を同じ距離送ると約\qty{1.3}{ms}、音なら約19分かかります。\\
極端な場合：屈折率が1（真空）なら $t=L/c$ で、いちばん短くなります。
\kotae 約\qty{2.0}{ms}（1000分の2秒）
\end{phreidai}

%% ================================================================
\bksubsection{5}{振動数のものさし}
\leavevmode\label{ph:波-0:sec5}この巻に出てくるゆれと波の振動数を、1本のものさしに並べてみましょう（図\ref{ph:波-0:fig-scale}）。ものさしの目もりは、1つ右へ進むごとに10倍になります（3-12 の対数の目もり）。
ビルのゆれ（約\qty{0.3}{Hz}）から、目に見える光（約 $5\times10^{14}$\,Hz）まで、振動数は15けたもちがいます。それでも、どれも $c=f\lambda$ で、振動数と波長と速さが結びついています。光の速さは $c=\qty{3.00e8}{m/s}$ なので、緑の光（波長\qty{550}{nm}）の振動数は $3.00\times10^8\div(550\times10^{-9})\approx5.45\times10^{14}$\,Hz です（波-13 ①）。
\begin{bkfigure}
\begin{tikzpicture}[x=1mm, y=1mm]
  % x = 8*(log10 f + 1)
  \draw[bkInk, line width=0.8pt, -{Stealth[length=2.4mm]}] (0,0) -- (131,0);
  \foreach \k in {0,...,16} {\draw[bkInk, line width=0.5pt] (\k*8,0) -- (\k*8,-1.6);}
  \foreach \x/\t in {8/1\,Hz, 32/1\,kHz, 56/1\,MHz, 80/1\,GHz, 104/$10^{12}$\,Hz, 128/$10^{15}$\,Hz} {\node[rxa lab, anchor=north] at (\x,-2.2) {\t};}
  % 帯（太い線）と点
  \fill[bkBlue] (4.2,0) circle (1.3);
  \draw[bkGreen!60!black, line width=3pt] (18.4,2.2) -- (42.4,2.2);
  \fill[bkOrange] (16.3,0) circle (1.3);
  \draw[bkPurple!80, line width=3pt] (53.8,2.2) -- (57.6,2.2);
  \draw[bkRed!70, line width=3pt] (59.2,-8.6) -- (64.0,-8.6);
  \draw[bkPurple!80, line width=3pt] (70.9,2.2) -- (71.8,2.2);
  \fill[bkPurple] (83.0,0) circle (1.3);
  \fill[bkRed] (122.3,0) circle (1.3);
  \draw[bkOrange!90, line width=3pt] (124.7,2.2) -- (127.2,2.2);
  % 字
  \node[rxa lab, align=center, anchor=south] at (6,3.5) {ビル\\0.3\,Hz};
  \node[rxa lab, align=center, anchor=north] at (16.3,-7.5) {A弦\\110\,Hz};
  \node[rxa lab, align=center, anchor=south] at (30.4,3.5) {人に聞こえる音\\20\,Hz〜20\,kHz};
  \node[rxa lab, align=center, anchor=south] at (55.7,3.5) {AM\\ラジオ};
  \node[rxa lab, align=center, anchor=north] at (61.6,-10.4) {超音波\\検査};
  \node[rxa lab, align=center, anchor=south] at (71.4,3.5) {FM\\ラジオ};
  \node[rxa lab, align=center, anchor=south] at (84.5,3.5) {Wi-Fi\\2.4\,GHz};
  \node[rxa lab, align=center, anchor=north] at (118,-7.5) {光ファイバの光\\1550\,nm};
  \node[rxa lab, align=center, anchor=south] at (122.5,3.5) {目に見える光\\380〜780\,nm};
\end{tikzpicture}
\bkfigcaption{振動数のものさし。目もりが1つ進むごとに10倍。太い線は範囲、丸は1つの値。左からゆれ（波-1〜4）、音（波-5〜11）、電波（波-14）、光（波-12・13）。}\label{ph:波-0:fig-scale}
\end{bkfigure}





%% ================================================================
\bksubsection{6}{4つの部と、14のなぜ？}
\leavevmode\label{ph:波-0:sec6}この巻は、序章と4つの部（14の章）でできています。
\begin{bkcenter}
\small
\begin{tblr}{colspec={Q[l,m,wd=19mm] X[l,m] Q[l,m,wd=31mm]}, width=\linewidth, rowsep=1.1pt, colsep=3pt,
  hlines={0.4pt,bkLine}, row{1}={bg=bkTheme!12, font=\sffamily\bfseries\footnotesize},
  row{2,7}={bg=bkTheme!6, font=\sffamily\bfseries\footnotesize}}
章 & 身近な「なぜ？」 & 答えのある所 \\
\SetCell[c=3]{l}{\phfaceicon[5mm]{1}\enspace 第1部　ゆれる：振動と制御} & & \\
\phchapref{波-1} & ブランコの1往復の時間は、何で決まる？ & ③ ふりこ \\
\phchapref{波-2} & 洗濯機は、なぜ回り始めにガタガタゆれることがある？ & ⑦ 不つり合い \\
\phchapref{波-3} & 高いビルの屋上の重りは、どうやってゆれをおさえる？ & ③⑥ 動吸振器 \\
\phchapref{波-4} & エアコンは、どうやって部屋をちょうどよい温度に保つ？ & ③ フィードバック \\
\SetCell[c=3]{l}{\phfaceicon[5mm]{2}\enspace 第2部　伝わる：波} & & \\
\phchapref{波-5} & スタジアムのウェーブでは、何が進んでいる？ & ① 波とは \\
\phchapref{波-6} & ギターの弦は、なぜおさえる場所で音の高さが変わる？ & ③ 弦の固有振動 \\
\phchapref{波-7} & 超音波の検査で、なぜおなかにゼリーをぬる？ & ④ はね返り \\
\phchapref{波-8} & 救急車のサイレンは、なぜ通りすぎると低く聞こえる？ & ⑤ ドップラー効果 \\
\end{tblr}
\end{bkcenter}
\begin{bkcenter}
\small
\begin{tblr}{colspec={Q[l,m,wd=19mm] X[l,m] Q[l,m,wd=31mm]}, width=\linewidth, rowsep=1.1pt, colsep=3pt,
  hlines={0.4pt,bkLine}, row{1,5}={bg=bkTheme!6, font=\sffamily\bfseries\footnotesize}}
\SetCell[c=3]{l}{\phfaceicon[5mm]{3}\enspace 第3部　はかる：音・振動の計測と信号} & & \\
\phchapref{波-9} & 雷が光ってから、なぜ少しおくれて音が聞こえる？ & ② 光と音の速さ \\
\phchapref{波-10} & コウモリは、暗やみでどうやってものの位置を知る？ & ① こだま \\
\phchapref{波-11} & 動画の車輪が、逆に回って見えることがあるのはなぜ？ & ② エイリアシング \\
\SetCell[c=3]{l}{\phfaceicon[5mm]{4}\enspace 第4部　曲がる：光} & & \\
\phchapref{波-12} & 水に入れたストローは、なぜ折れて見える？ & ③ 屈折の法則 \\
\phchapref{波-13} & 空は、なぜ青い？　夕焼けは、なぜ赤い？ & ①⑧ 散乱 \\
\phchapref{波-14} & スマホの声は、どうやって電波にのせて送る？ & ⑨ 声が電波になるまで \\
\end{tblr}
\bktabcaption{この巻の地図。「答えのある所」の丸数字は、その章の「物理の道具」の節の番号。}\label{ph:波-0:tab-chizu}
\end{bkcenter}
\begin{itemize}[leftmargin=1.2\zw, itemsep=2pt, topsep=2pt, label={\color{bkTheme}●}]
\item \textbf{第1部　ゆれる}：ばねとおもり、ふりこ、ビル。1か所でくり返すゆれと、それを止める・ねらいどおりに動かすしくみ（波-1〜4）。
\item \textbf{第2部　伝わる}：ゆれが、ひも・空気・水・地面を伝わる波。はね返り、重なり、エネルギーの運び方（波-5〜8）。
\item \textbf{第3部　はかる}：音の大きさ、超音波、振動のセンサー、音を数字にする方法（波-9〜11）。
\item \textbf{第4部　曲がる}：光の反射・屈折・レンズ、光の波としての性質、波で情報を送るしくみ（波-12〜14）。
\end{itemize}
表\ref{ph:波-0:tab-chizu}は、それぞれの章で答える身近な「なぜ？」と、その答えがある所です。章の札の横の数は、その章が始まるページです。どの章から読んでもよいのですが、前の章の道具を使うので、はじめは順に読むのがおすすめです。序章の\bkhatenaref の答えは、この序章の最後（⑨）にあります。身近な疑問から、数学と物理のどの章を順に読めばよいかを示した「なぜ？の道すじ」も、巻末にあります。

%% ================================================================
\bksubsection{7}{数学の本とのつながり}
\leavevmode\label{ph:波-0:sec7}ゆれと波の式の多くは、数学の本でもう計算しています。たとえば、ばねとおもりの式は 6-6、ゆすり続けたときの共振は 6-13、弦の波の式は 8-13 で解きました。この巻では、同じ計算をくり返すかわりに、「（6-13⑩）」のように数学の章と節を指します。そして、物理の本ならではの、次の3つを中心に書きます。
\begin{itemize}[leftmargin=1.2\zw, itemsep=2pt, topsep=2pt, label={\color{bkTheme}●}]
\item その式は、何を無視して（どんな仮定で）作ったのか。
\item 本物の物（ビル・楽器・光ファイバ）の値を入れると、何が起きるのか。
\item 仮定がくずれると、何が変わるのか。
\end{itemize}
\begin{bkcenter}
\small
\begin{tblr}{colspec={Q[l,m,wd=46mm] X[l,m] Q[l,m,wd=28mm]}, width=\linewidth, rowsep=1.1pt, colsep=3pt,
  hlines={0.4pt,bkLine}, row{1}={bg=bkTheme!12, font=\sffamily\bfseries\footnotesize}}
数学の本の章 & そこで解いたこと & この巻の章 \\
6-6・6-7・6-12 & ばねとおもり・ふりこ、ゆれが弱まる & 波-1 \\
6-13 & ゆすり続けたときの共振、Q 値 & 波-2 \\
6-14・7-8・7-9 & 2つのおもりのゆれ、固有値 & 波-3 \\
9-8・9-9・9-10 & 伝達関数、安定、状態方程式 & 波-4 \\
8-13 & 波動方程式、形を保って進む波 & 波-5 \\
\end{tblr}
\end{bkcenter}
\begin{bkcenter}
\small
\begin{tblr}{colspec={Q[l,m,wd=46mm] X[l,m] Q[l,m,wd=28mm]}, width=\linewidth, rowsep=1.1pt, colsep=3pt,
  hlines={0.4pt,bkLine}}
9-1・9-2・9-3 & 弦の倍音、フーリエ級数 & 波-6 \\
4-9 & 波のたし合わせ、うなり & 波-8・波-13 \\
3-12 & 対数とデシベル & 波-9 \\
9-4 & スペクトル、サンプリング & 波-11 \\
7-3 & 光線の行列 & 波-12 \\
9-12 & 音楽の圧縮、Wi-Fi & 波-14 \\
\end{tblr}
\bktabcaption{数学の本で解いた計算と、それを使うこの巻の章。}\label{ph:波-0:tab-math}
\end{bkcenter}
\noindent 数学の本をまだ読んでいなくても、だいじょうぶです。★が5つの内容をふくむ章は、「目でわかる」「手でわかる」「式でわかる」の3段に分けてあり、はじめの2段は、図と表計算を中心に、★3つまでの道具でたどれるように書いてあります。

%% ================================================================
\Needspace{8\baselineskip}
\bksubsection{8}{この本の読み方}
\leavevmode\label{ph:波-0:sec8}
\begin{itemize}[leftmargin=1.2\zw, itemsep=2pt, topsep=2pt, label={\color{bkTheme}●}]
\item \textbf{★の数}：内容の高さの目安です。★1つは「はじめの一歩」、★2つは「ふもとを歩く」、★3つは「山を登る」、★4つは「尾根を歩く」、★5つは「頂上をめざす」。物理の本の★は、物理の内容の高さで付けてあり、数学の本の★とは別の物差しです。
\item \textbf{章の札}：「8-13」のような札は数学の本の章、「波-5」のような札はこの本の章、「力-2」「熱-4」は第1巻『力学』・第2巻『熱学』の章です。札の色は、その本の色です。本文で「（8-13⑩）」と書いたら、数学の本の8-13の⑩の節を見てください。
\item \textbf{物理の本の囲み}：「この式の読み方」「身近な例では」「〇〇を2倍にすると？」「まだわかっていないこと」「この先の学び」「一般の式から、よく使う式へ」などがあります。見本は、前付の「この本の使い方」にあります。
\item \textbf{本物の値}：音の速さや屈折率など、この巻で使う物の値は、巻末の「物性値・定数・単位」にまとめてあります。
\end{itemize}
\begin{bkfigure}
\begin{tikzpicture}[x=1mm, y=1mm]
  \rxatalk{0}{sora}{ソラ}{音って、空気が耳まで飛んでくるのかと思っていたよ。}
  \rxatalk{-20}{mina}{ミナ}{わたしは、ビルがゆれる話と光の話が、同じ本にあるのがふしぎ……。}
  \rxatalk{-44}{delta}{デルタ}{空気は、その場でゆれるだけで、進むのは「ゆれ」という形だよ（波-5）。ゆれも光も、同じ $\sin$ と同じ $c=f\lambda$ で書けるから、1冊の本になるんだ。★が5つの所は3段に分けてあるから、「ここまでで十分」の旗まで読めば、その章のゴールだよ。}
\end{tikzpicture}
\end{bkfigure}
\begin{taisetsu}
\def\tsl#1{\par\noindent\hangindent=6.2\zw\hangafter=1\makebox[6.2\zw][l]{\sffamily\bfseries\small\color{bkDeep}#1}}
\tsl{周期と振動数}周期 $T$〔s〕は1往復の時間、振動数 $f$〔Hz〕は1秒の往復の回数。$f=1/T$。
\tsl{波の速さ}$c=f\lambda$。この巻では波の速さを $c$ と書く（高校の教科書の $v$・$V$）。
\tsl{$c$ の意味}第1部では減衰係数（ゆれを弱める力の強さ）、第2・3部では波の速さ、第4部では光の速さ。同じ章で2つの意味には使わない。
\tsl{音の速さ}\qty{20}{\celsius}の空気で\qty{343}{m/s}。中学理科とくらべる例題では\qty{340}{m/s}。
\tsl{答えのけた}答えは、ふつう2けた。数学の本と同じ数を使う例題などは3けたで書く。
\end{taisetsu}
\noindent 記号の読み方と意味は、巻末の「記号カード」にまとめてあります。

%% ================================================================
\Needspace{14\baselineskip}
\bksubsection{9}{はてな？の答え：ゆれも、音も、光も波}
\leavevmode\label{ph:波-0:sec9}\bkhatenaref に答えましょう。
\begin{bkfigure}
\begin{tikzpicture}[x=1mm, y=1mm]
  \tikzset{tag/.style={draw=#1, line width=0.9pt, fill=white, rounded corners=1.6mm, text width=58mm, align=left, inner sep=1.8mm, font=\sffamily\footnotesize}}
  \node[tag=bkBlue, anchor=north west] (n1) at (0,0) {\phfaceicon[5mm]{1}\enspace\textbf{ゆれる}\\ビルは、自分のリズム（固有振動数）でくり返しゆれる。\\\phchapref{波-1}\enspace\phchapref{波-3}};
  \node[tag=bkOrange!80!black, anchor=north west] (n2) at (66,0) {\phfaceicon[5mm]{2}\enspace\textbf{伝わる}\\弦のゆれが空気を伝わって、音になる。\\\phchapref{波-5}\enspace\phchapref{波-6}};
  \node[tag=bkGreen!70!black, anchor=north west] (n3) at (0,-31) {\phfaceicon[5mm]{3}\enspace\textbf{はかる}\\音や地面のゆれを測って、数にする。\\\phchapref{波-9}\enspace\phchapref{波-11}};
  \node[tag=bkPurple, anchor=north west] (n4) at (66,-31) {\phfaceicon[5mm]{4}\enspace\textbf{曲がる}\\光も波。曲がり、重なり、情報を運ぶ。\\\phchapref{波-12}\enspace\phchapref{波-14}};
\end{tikzpicture}
\bkfigcaption{はてな？の答え。ビルのゆれ・楽器の音・光は、どれも「くり返す」と「伝わる」の物理で読める。}
\end{bkfigure}
\begin{enumerate}[label={\bkcircled[bkBlue]{\arabic*}}, leftmargin=2.4\zw, itemsep=2pt, topsep=2pt]
\item \textbf{どれも、くり返す}　ビルは約3秒ごとに（\bkreidairef{2}）、ギターの弦は1秒に110回、緑の光は1秒に約 $5\times10^{14}$ 回、同じ動きをくり返します。くり返しは、周期 $T$ と振動数 $f$ で表せます（②）。
\item \textbf{どれも、同じ形の式}　くり返すゆれは、時間の $\sin$ の形で書けます。伝わる波は、時間と場所の $\sin$ の形で書けます（図\ref{ph:波-0:fig-same}）。どちらも、もとは運動方程式（力-2）から出てくる式です（波-1・波-5）。
\item \textbf{どれも、$c=f\lambda$}　音も光も、伝わる波は、1周期の間に1波長進みます（②）。ちがうのは、振動数の大きさ（15けたもちがう：図\ref{ph:波-0:fig-scale}）と、何がゆれているか（ビルの柱、空気、光では電場と磁場）です。
\end{enumerate}
だから、ビルのゆれを止める考えと、楽器の音を作る考えと、光で情報を送る考えを、1冊の本で、同じ道具を使って学べるのです。
\begin{phketa}[ゆれと波の数]
\phketaline{30階建てのビルの周期}{約\qty{3}{s}（約\qty{0.3}{Hz}）}
\phketaline{ギターの A弦}{\qty{110}{Hz}、長さ\qty{0.65}{m}}
\phketaline{人に聞こえる音}{約\qty{20}{Hz}〜\qty{20}{kHz}}
\phketaline{空気の中の音の速さ（\qty{20}{\celsius}）}{\qty{343}{m/s}}
\phketaline{地震の P波・S波}{秒速 約7 km・約4 km}
\phketaline{真空の光の速さ}{\qty{3.00e8}{m/s}（ガラスの中で約\qty{2.0e8}{m/s}）}
\phboxhead{答えが桁で合っているか}
音の速さは光の約100万分の1です。雷までの距離が何千 km になったり、光ファイバで数分かかったりしたら、km と m、ms と s を取り違えています。Hz と「周期の秒」は逆数なので、ビルの振動数が何百 Hz になったら、$f$ と $T$ を取り違えています。
\end{phketa}



\begin{column}{音は真空を伝わらない。光は伝わる}
音は、空気の分子が押し合って伝わる波なので、空気のない真空の中は伝わりません。宇宙では、となりの宇宙船の爆発の音は聞こえないのです。一方、太陽の光は、ほとんど真空の宇宙を通って、地球まで届きます。光は、ゆれる電場と磁場の波（電磁波）で、ゆれを伝える物（媒質）がなくても進めるからです（物理シリーズの『電磁気学』（電-15））。それでも、音と光は、どちらも $c=f\lambda$ にしたがい、はね返り・曲がり・重なります。何がゆれているかはちがっても、波としてのふるまいは同じなのです。
\end{column}

\begin{kotoba}\small
\bkword{振動（ゆれ）}{1か所で、行ったり来たりをくり返す動き。}
\bkword{波}{ゆれが、周りへ伝わっていくこと。物そのものは遠くへ運ばれない。}
\bkword{周期・振動数}{1往復の時間 $T$〔s〕と、1秒の往復の回数 $f$〔Hz〕。$f=1/T$。}
\bkword{波長}{波の形がくり返す長さ $\lambda$〔m〕。}
\bkword{固有振動数}{物が、自分のリズムでゆれるときの振動数。}
\end{kotoba}

\begin{furikaeri}
\begin{checklist}
\item ゆれの周期と振動数、波の波長と速さを、身近な例で説明できる。\bkfurref{ph:波-0:sec2}
\item 何を無視したかを書いてから、波が届くまでの時間や距離を、7つの段で求められる。\bkfurref{ph:波-0:ex3}
\item 身近な「なぜ？」の答えが、この巻のどの章にあるかを探せる。\bkfurref{ph:波-0:sec6}
\end{checklist}
\end{furikaeri}
\bkfukidashi{sora}{ビルのゆれも、ギターの音も、光も、同じ $\sin$ の形なんだね。まずは波-1 で、ばねとおもりのゆれから始めよう。}
